% Exam Template for UMTYMP and Math Department courses
%
% Using Philip Hirschhorn's exam.cls: http://www-math.mit.edu/~psh/#ExamCls
%
% run pdflatex on a finished exam at least three times to do the grading table on front page.
%
%%%%%%%%%%%%%%%%%%%%%%%%%%%%%%%%%%%%%%%%%%%%%%%%%%%%%%%%%%%%%%%%%%%%%%%%%%%%%%%%%%%%%%%%%%%%%%

% These lines can probably stay unchanged, although you can remove the last
% two packages if you're not making pictures with tikz.
\documentclass[11pt]{exam}
\RequirePackage{amssymb, amsfonts, amsmath, latexsym, verbatim, xspace, setspace, makecell, rotating, tcolorbox}
\RequirePackage{tikz, pgflibraryplotmarks}

% By default LaTeX uses large margins.  This doesn't work well on exams; problems
% end up in the "middle" of the page, reducing the amount of space for students
% to work on them.
\usepackage[margin=1in]{geometry}


% Here's where you edit the Class, Exam, Date, etc.
\newcommand{\class}{Math 20250 (Practice)}
\newcommand{\term}{Autumn 2025}
\newcommand{\examnum}{Final Exam}
\newcommand{\examdate}{12/05/25}
\newcommand{\timelimit}{120 Minutes}

\newcommand\aug{\fboxsep=-\fboxrule\!\!\!\fbox{\strut}\!\!\!}

\newcommand{\Span}{\mathrm{Span}}


% For an exam, single spacing is most appropriate
\singlespacing
% \onehalfspacing
% \doublespacing

% For an exam, we generally want to turn off paragraph indentation
\parindent 0ex

\begin{document} 

% These commands set up the running header on the top of the exam pages
\pagestyle{head}
\firstpageheader{}{}{}
\runningheader{Math 20250}{\examnum\ - Page \thepage\ of \numpages}{\examdate}
\runningheadrule

\begin{flushright}
\begin{tabular}{p{2.6in} r l}
\textbf{\class} \\
\textbf{\term} &&\\
\textbf{\examnum} &&\\
\textbf{\examdate} &&\\
\textbf{Time Limit: \timelimit} & \textbf{Name (Print):} \makebox[2.5in]{\hrulefill}
\end{tabular}\\
\end{flushright}
\rule[1ex]{\textwidth}{.1pt}

\bigskip

This exam contains \numpages\ pages (including this cover page) and three problems.\\

This exam is out of a total of 40 possible points.\\

Check to see if any pages are missing.  Enter
all requested information on the top of this page, and \textbf{put your initials
on the top of every page, in case the pages become separated.}\\

You may \emph{not} use your books, notes, a computer, or a calculator on this exam.\\

\textbf{Unless specified otherwise, you are required to show your work on each problem on this exam.}  The following rules apply:\\[-5pt]

\begin{minipage}[t]{6.25in}
\vspace{0pt}
\begin{itemize}

%\item \textbf{If you use a ``fundamental theorem'' you must indicate this and explain
%why the theorem may be applied.}

\item \textbf{Organize your work}, in a reasonably neat and coherent way, in
the space provided. Work that is scattered all over the page without a clear ordering will 
receive very little credit.  

\item \textbf{On problems that require work to be shown, answers that are mysterious or unsupported will not receive full
credit.}  A correct answer, unsupported by calculations, explanation,
or algebraic work will receive no credit; an incorrect answer supported
by substantially correct calculations and explanations might still receive
partial credit.


\item \textbf{Please try to contain your answer in the space provided.} If you need more space to work on a problem, use the back of the pages; please clearly indicate when you have done this.
\end{itemize}

%Do not write in the table to the right.


\end{minipage}

\bigskip
\bigskip
\medskip

Following the University of Chicago's policy on academic honesty, I pledge that the work shown on this exam is entirely my own and that I have followed the rules stated above. 

\bigskip
\bigskip
\bigskip

\qquad Signature: \hrulefill

\hspace*{0mm}\phantom{}

\vfill


%\hfill
%\begin{minipage}[t]{2.3in}
%\vspace{0pt}
%\cellwidth{3em}
%\gradetablestretch{2}
%\vqword{Problem}
%\addpoints % required here by exam.cls, even though questions haven't started yet.	
%\gradetable[v]%[pages]  % Use [pages] to have grading table by page instead of question

%\end{minipage}

\newpage % End of cover page

%%%%%%%%%%%%%%%%%%%%%%%%%%%%%%%%%%%%%%%%%%%%%%%%%%%%%%%%%%%%%%%%%%%%%%%%%%%%%%%%%%%%%
%
% See http://www-math.mit.edu/~psh/#ExamCls for full documentation, but the questions
% below give an idea of how to write questions [with parts] and have the points
% tracked automatically on the cover page.
%
%
%%%%%%%%%%%%%%%%%%%%%%%%%%%%%%%%%%%%%%%%%%%%%%%%%%%%%%%%%%%%%%%%%%%%%%%%%%%%%%%%%%%%%

{\bf Problem 1.} (10 points) 
\iffalse
Recall that a linear transformation $f\colon V \to W$ is an \emph{isomorphism} if and only if there exists a linear transformation $g\colon W\to V$ such that the composites
$$g\circ f = \mathrm{id}_{V}, \;\;\; f\circ g =\mathrm{id}_W $$
are the identity maps.

\bigskip
\fi
Let $V,W$ be finite dimensional vector spaces over a field $\mathbb{F}$. Show that there is an \emph{injective} linear transformation $f\colon V \to W$ if and only if $\dim_{\mathbb{F}}V\leq \dim_{\mathbb{F}}W$.

\newpage


{\bf Problem 2.} (10 points) Consider the following system of linear equations
$$\begin{cases} x +y - z=1,\\ x+z=-3,\\ -x+y+2z=1.\end{cases}$$
Solve for $x,y,z$ assuming that $x,y,z \in \mathbb{F}_5$ are elements of the finite field $\mathbb{F}_5$. Recall that $\mathbb{F}_5=\{0,1,2,\dots, 4\}$ where addition, subtraction, and multiplication are defined according to the following rule: first add/subtract/multiply the integers $a,b\in\mathbb{F}_5$ together, and then take the remainder after division by $5$. 

\noindent For example, $9=4$ and $2\cdot 3=1$ in the field $\mathbb{F}_5$.
\iffalse
\bigskip

{\bf Solution. }

\begin{align*}
\begin{bmatrix}
1 & 1 & -1 &\aug& 1\\
1 & 0 & 1 &\aug& -3\\
-1 & 1 & 2 &\aug& 1
\end{bmatrix}
&\xrightarrow[R_3\to R_3+R_1]{R_2\to R_2-R_1}
\begin{bmatrix}
1 & 1 & -1 &\aug& 1\\
0 & -1 & 2 & \aug& -4\\
0 & 2 & 1 &\aug& 2
\end{bmatrix}\\[5pt]
&\xrightarrow{R_2\to -R_2}
\begin{bmatrix}
1 & 1 & -1 &\aug& 1\\
0 & 1 & -2 &\aug& 4\\
0 & 2 & 1 &\aug& 2
\end{bmatrix}\\[5pt]
&\xrightarrow{R_3\to R_3-2R_2}
\begin{bmatrix}
1 & 1 & -1 &\aug& 1\\
0 & 1 & -2 &\aug& 4\\
0 & 0 & 5 &\aug& -6
\end{bmatrix}.
\end{align*}
Since $5=0$ in the field $\mathbb{F}_5$, the last equation has the form $0\cdot z = -6=4$. This equation does not have any solutions, so the linear system has no solutions.
\fi
\newpage

{\bf Problem 3.} Let $T\colon \mathbb{C}^3 \to \mathbb{C}^3$ be a linear transformation given the matrix
$$A=\begin{bmatrix} -2 & 2 & 6 \\ 5 & 1 & -6 \\ -5 &2 &10 \end{bmatrix} \in M_{3\times 3}(\mathbb{C}). $$
\begin{enumerate}
\item (3 points) Find the characteristic polynomial $p_{A}(t)$ of the matrix $A$ and the eigenvalues of the linear transformation $T$.
\iffalse
\bigskip

{\bf Solution.} By the definition, $f_A(t)=\det(tI_3-A)$. We have
$$f_A(t)=\det(tI_3-A)=\det
\begin{bmatrix}
t +2 & -2 & -6\\
-5 &t-1 & 6 \\
5 & -2 & t-10
\end{bmatrix}.
$$
We compute this determinant by the Laplace expansion across the first column:
\begin{align*}
\det
\begin{bmatrix}
t +2 & -2 & -6\\
-5 &t-1 & 6 \\
5 & -2 & t-10
\end{bmatrix}&=
(t+2)\cdot \det 
\begin{bmatrix}
t-1 & 6\\
-2 & t-10
\end{bmatrix} + 5\cdot \det 
\begin{bmatrix}
-2 & -6\\
-2 & t-10
\end{bmatrix} \\
&+5\cdot \det 
\begin{bmatrix}
-2 & -6\\
t-1 & 6
\end{bmatrix} \\
&=(t+2)((t-1)(t-10)+12)+5(-2(t-10)-12)\\
&+5(-12+6(t-1))\\
&=t^3-9t^2+20t-6.
\end{align*}
Therefore, $f_A(t)=t^3-9t^2+20t-6$.
\medskip

The eigenvalues of $T$ are the roots of the characteristic polynomial $f_A(t)$. By substitution, we observe that $\lambda_1=3$ is a root of $f_A(t)$. By the long division, we get that
$$f_A(t)= t^3-9t^2+20t-6 = (t-3)(t^2-6t+2).$$
By the quadratic formula, that two other roots are $\lambda_2=3-\sqrt{7}$ and $\lambda_3=3+\sqrt{7}$. Therefore, the eigenvalues of $T$ are $\lambda_1 = 3$, $\lambda_2=3-\sqrt{7}$, and $\lambda_3=3+\sqrt{7}$
\fi
\newpage
\item (3 points) Find the eigenvectors of the linear transformation $T$. Is $T$ diagonalizable?

\iffalse
\bigskip

{\bf Solution.} Since the linear transformation $T$ has three distinct eigenvalues, $T$ is diagonalizable. We find the eigenvectors.
\begin{itemize}
\item $\lambda_1=3$. Then $\ker(A-\lambda_1I_3)=\ker\begin{bmatrix} -5 & 2 & 6 \\ 5 & -2 & -6 \\ -5 &2 &7 \end{bmatrix} = \Span_{\mathbb{C}}\begin{bmatrix} 2 \\ 5 \\ 0 \end{bmatrix} $. So, the eigenvector $\mathbf{v}_1$ corresponding to the eigenvalue $\lambda_1=3$ is $\mathbf{v}_1=\begin{bmatrix} 2 \\ 5 \\ 0 \end{bmatrix}$.

\item $\lambda_2=3-\sqrt{7}$. Then $\ker(A-\lambda_2I_3)=\ker\begin{bmatrix} -5+\sqrt{7} & 2 & 6 \\ 5 & -2+\sqrt{7} & -6 \\ -5 &2 & 7+\sqrt{7} \end{bmatrix}$. In order to find the kernel, we find the reduced echelon form of the matrix:
\begin{align*}
\begin{bmatrix} -5+\sqrt{7} & 2 & 6 \\
5 & -2+\sqrt{7} & -6 \\
-5 &2 &7+\sqrt{7} 
\end{bmatrix}
&\xrightarrow{R_2\to R_2+R_1}
\begin{bmatrix} -5+\sqrt{7} & 2 & 6 \\
\sqrt{7} & \sqrt{7} & 0 \\
-5 &2 &7+\sqrt{7}
\end{bmatrix}
\\[5pt]
&\xrightarrow{R_2\to \frac{1}{\sqrt{7}}R_2}
\begin{bmatrix} -5+\sqrt{7} & 2 & 6 \\
1 & 1 & 0 \\
-5 &2 &7+\sqrt{7}
\end{bmatrix}
\\[5pt]
&\xrightarrow{R_2 \leftrightarrow R_1}
\begin{bmatrix}
1 & 1 & 0 \\
 -5+\sqrt{7} & 2 & 6 \\
-5 &2 &7+\sqrt{7}
\end{bmatrix}
\\[5pt]
&\xrightarrow[R_3\to R_3+5R_1]{R_2 \to R_2-(-5+\sqrt{7})R_1}
\begin{bmatrix}
1 & 1 & 0 \\
0 & 7-\sqrt{7} & 6 \\
0 & 7 &7+\sqrt{7}
\end{bmatrix}
\\[5pt]
&\xrightarrow{R_3 \to \frac{1}{7}R_3}
\begin{bmatrix}
1 & 1 & 0 \\
0 & 7-\sqrt{7} & 6 \\
0 & 1 &\frac{7+\sqrt{7}}{7}
\end{bmatrix}
\\[5pt]
&\xrightarrow{R_2 \leftrightarrow R_3}
\begin{bmatrix}
1 & 1 & 0 \\
0 & 1 &\frac{7+\sqrt{7}}{7} \\
0 & 7-\sqrt{7} & 6 \\
\end{bmatrix}
\\[5pt]
&\xrightarrow{R_3 \to R_3- (7-\sqrt{7})R_2}
\begin{bmatrix}
1 & 1 & 0 \\
0 & 1 &\frac{7+\sqrt{7}}{7} \\
0 & 0 & 6 - \frac{(7+\sqrt{7})(7-\sqrt{7})}{7}
\end{bmatrix}
\\[5pt]
&=
\begin{bmatrix}
1 & 1 & 0 \\
0 & 1 &\frac{7+\sqrt{7}}{7} \\
0 & 0 & 0
\end{bmatrix}.
\end{align*}
Therefore,  $\ker(A-\lambda_2I_3)=\Span_{\mathbb{C}}\begin{bmatrix} \frac{1}{7}(7+\sqrt{7}) \\  -\frac{1}{7}(7+\sqrt{7}) \\ 1 \end{bmatrix} $. So, the eigenvector $\mathbf{v}_2$ corresponding to the eigenvalue $\lambda_2=3-\sqrt{7}$ is $\mathbf{v}_2=\begin{bmatrix} \frac{1}{7}(7+\sqrt{7}) \\  -\frac{1}{7}(7+\sqrt{7}) \\ 1 \end{bmatrix}$.

\item $\lambda_3=3+\sqrt{7}$. Then $\ker(A-\lambda_3I_3)=\ker\begin{bmatrix} -5-\sqrt{7} & 2 & 6 \\ 5 & -2-\sqrt{7} & -6 \\ -5 &2 & 7-\sqrt{7} \end{bmatrix}$. By replacing $\sqrt{7}$ with $-\sqrt{7}$ in the previous argument, we find $\ker(A-\lambda_3I_3)=\Span_{\mathbb{C}}\begin{bmatrix} \frac{1}{7}(7-\sqrt{7}) \\  -\frac{1}{7}(7-\sqrt{7}) \\ 1 \end{bmatrix} $. So, the eigenvector $\mathbf{v}_3$ corresponding to the eigenvalue $\lambda_3=3+\sqrt{7}$ is $\mathbf{v}_3=\begin{bmatrix} \frac{1}{7}(7-\sqrt{7}) \\  -\frac{1}{7}(7-\sqrt{7}) \\ 1 \end{bmatrix}$.
\end{itemize}
\fi
\newpage
\item (4 points) Show that there are finitely many matrices $B\in M_{3\times 3}(\mathbb{C})$ such that $B^2=A$.

\iffalse
\bigskip

{\bf Solution.} By part (a), the matrix $A$ has the form $A=CDC^{-1}$, where $C$ is an invertible matrix and 
$$D= \begin{bmatrix} 3 & 0 & 0 \\
0 & 3-\sqrt{7} & 0 \\
0 & 0 &3+\sqrt{7} 
\end{bmatrix}.
$$
Therefore, any square root matrix $B$ has the form $B=CD_1C^{-1}$, where $C$ is the same invertible matrix and $D_1^2=D$. Since $D$ is a diagonal matrix with pairwise distinct entries, the matrix $D_1$ is again diagonal of the form
$$D_1= \begin{bmatrix} \pm \sqrt{3} & 0 & 0 \\
0 & \pm \sqrt{3-\sqrt{7}} & 0 \\
0 & 0 & \pm \sqrt{3+\sqrt{7}}
\end{bmatrix}.
$$
In total, there are $2^3$ possibilities for different signs, so there are exactly $8$ possible square root matrices.

\fi
\newpage
\end{enumerate}

{\bf Problem 4.} (10 points) Let $A \in M_{n\times n}(\mathbb{C})$ be an $n \times n$ complex matrix. Suppose that all entries of $A$ are real numbers. 

\medskip

Suppose that $\lambda \in \mathbb{C}$ is an eigenvalue of $A$. Show that its complex conjugate $\bar{\lambda}\in \mathbb{C}$ is also an eigenvalue.


\end{document}
